\documentclass[12pt,a4paper]{article}
\usepackage[utf8]{inputenc}
\usepackage{amsmath,amssymb,amsthm}
\usepackage{algorithm}
\usepackage{algpseudocode}
\usepackage{geometry}
\geometry{margin=1in}

\newtheorem{theorem}{Theorem}
\newtheorem{lemma}[theorem]{Lemma}
\newtheorem{corollary}[theorem]{Corollary}
\theoremstyle{definition}
\newtheorem{definition}[theorem]{Definition}

\title{Problem 74: Digit Factorial Chains}
\author{Project Euler Solution}
\date{}

\begin{document}
\maketitle

\section{Problem Statement}

Starting with any positive integer and repeatedly replacing it with the sum of the factorials of its digits, a chain is formed that eventually cycles. How many chains, with starting number below $10^6$, contain exactly 60 non-repeating terms?

\section{Mathematical Analysis}

\begin{definition}[Digit Factorial Function]
The function $f: \mathbb{N} \to \mathbb{N}$ is defined by $f(n) = \sum_{i=0}^{k} d_i!$ where $d_0 d_1 \cdots d_k$ is the decimal representation of $n$.
\end{definition}

\begin{definition}[Chain Length]
The chain length $L(n)$ is the cardinality of the set $\{n, f(n), f^2(n), \ldots\}$ up to (but not including) the first repeated element.
\end{definition}

\begin{theorem}[Eventual Periodicity]
For all $n \geq 1$, the orbit $(n, f(n), f^2(n), \ldots)$ eventually enters a cycle. For $n < 10^7$, we have $f(n) \leq 2\,540\,160$.
\end{theorem}

\begin{proof}
A $k$-digit number $n$ satisfies $n \geq 10^{k-1}$ and $f(n) \leq 362880k$. For $k \geq 8$, $10^{k-1} > 362880k$, so $f(n) < n$. Thus $f$ eventually maps any input into $[1, 2\,540\,160]$, a finite set, guaranteeing periodicity.
\end{proof}

\begin{theorem}[Complete Cycle Classification]
The function $f$ has exactly these cycles:
\begin{itemize}
    \item Fixed points: $\{1\}$, $\{2\}$, $\{145\}$, $\{40585\}$.
    \item 2-cycle: $\{871, 45361\}$.
    \item 3-cycle: $\{169, 363601, 1454\}$.
\end{itemize}
\end{theorem}

\begin{proof}
Every cycle element lies in $[1, 2\,540\,160]$ by Theorem~1. Exhaustive iteration from every starting point in this range confirms exactly the listed cycles. The fixed points are verified by direct computation: $f(1) = 1$, $f(2) = 2$, $f(145) = 145$, $f(40585) = 40585$. The 2-cycle: $f(871) = 45361$, $f(45361) = 871$. The 3-cycle: $f(169) = 363601$, $f(363601) = 1454$, $f(1454) = 169$.
\end{proof}

\begin{lemma}[Chain Length Recursion]
If $f(n) = m$ and $n$ is not part of a cycle, then $L(n) = 1 + L(m)$.
\end{lemma}

\begin{proof}
The chain starting at $n$ is $(n, m, f(m), \ldots)$. Since $n$ is not in a cycle, it does not reappear, so $L(n) = 1 + L(m)$.
\end{proof}

\begin{corollary}[Memoization]
Once $L(m)$ is known, the chain length of any non-cyclic predecessor of $m$ is immediately determined. This enables dynamic programming over the finite state space $[1, 2\,540\,160]$.
\end{corollary}

\section{Algorithm}

\begin{algorithm}[H]
\caption{Count chains of length 60}
\begin{algorithmic}[1]
\Require $N = 10^6$
\State Precompute $\mathrm{fact}[d] = d!$ for $d = 0, \ldots, 9$
\State Initialize memoization cache $L[\cdot]$
\State $\mathrm{count} \gets 0$
\For{$n = 1$ \textbf{to} $N - 1$}
    \State Trace chain from $n$, recording visited values
    \State On hitting a cached value or detecting a cycle, fill cache via Lemma~3
    \If{$L[n] = 60$}
        \State $\mathrm{count} \gets \mathrm{count} + 1$
    \EndIf
\EndFor
\State \Return count
\end{algorithmic}
\end{algorithm}

\section{Complexity}

\textbf{Time:} Each value in $[1, M]$ ($M = 2\,540\,160$) is processed at most once. Computing $f$ costs $O(\log n)$ per step. Total: $O((N + M) \log N) = O(N \log N)$.

\textbf{Space:} $O(N + M) = O(N)$ for the cache.

\section{Answer}

\[
\boxed{402}
\]

\end{document}
